\podsumowanie

% Odpowiednik sekcji Conclusion i Future Work. Objętość: 2-3 strony.
% Podsumowanie nawiązuje do Wprowadzenia: wraca do celu i pytania.
% Nie wprowadzaj tu niczego nowego i nie powołuj się na nowe źródła.

\wyroznienie{Realizacja celu.} Wróć do celu postawionego we Wprowadzeniu i napisz wprost, w jakim stopniu został osiągnięty.

\wyroznienie{Odpowiedź na pytanie.} Odpowiedz na postawione pytanie badawcze. Odpowiedź ma wynikać z tego, co pokazałeś w rozdziałach drugim i trzecim.

\wyroznienie{Wnioski.} Trzy do pięciu wniosków, każdy w osobnym akapicie, uporządkowanych od najważniejszego.

\wyroznienie{Ograniczenia.} Czego narzędzie nie potrafi i przy jakich danych zawodzi.

\wyroznienie{Perspektywy rozwoju.} Konkretne kierunki dalszej pracy nad aplikacją. Wymóg ze Standardów prac dyplomowych EPI — nie pomijaj tej części i nie zbywaj jej jednym zdaniem.
